\chapter{Conclusion}\label{chap:concl}

This chapter summarises the work carried out, evaluates the extent to which the
original objectives were achieved, reflects on the technical decisions made during
the project, and acknowledges the limitations of the delivered system.

% ─────────────────────────────────────────────────────────────────────────────
\section{Summary of Work}\label{sec:concl_summary}
% ─────────────────────────────────────────────────────────────────────────────

Task~Flow was conceived in response to a concrete problem: existing task management
tools are web-first services that become inaccessible or functionally degraded when
network connectivity is lost, yet mobile or remote workers routinely face exactly those
conditions.
The objective was to build a self-contained desktop application that delivers the full
feature set of a modern task management platform while remaining fully operational
without an internet connection and automatically synchronising data when connectivity
is restored.

The six objectives established in Section~\ref{sec:intro_aims} were all met:

\begin{enumerate}

  \item \textbf{Cross-platform desktop application.}
        A Tauri~v2 shell embeds the ASP.NET Core backend and the React~19 frontend
        in a single redistributable package.
        The application packages into a Windows installer of approximately
        35~MB with a cold-start time of approximately 1.5~seconds.
        Installers were produced for Windows (NSIS and MSI).
        No prior installation of Node.js, .NET, Rust, or any other runtime is required
        on the target machine (Section~\ref{sec:res_distribution}).

  \item \textbf{Offline-first task management.}
        Every write operation produces both an immediate SQLite record and a
        \texttt{SyncOutboxEntry} in the outbox queue.
        The \texttt{OfflineSyncService} replays the queue when connectivity is restored
        and \texttt{BulkSyncStartupService} handles the case in which the user was offline
        during the preceding session.
        Twenty-entry offline-then-sync scenarios were verified without data loss
        (Section~\ref{sec:res_sync}).

  \item \textbf{Multiple task visualisation modes.}
        Five views --- Default, Kanban, Table, Gantt, and Calendar --- each present the
        task dataset through a different organisational lens.
        The Kanban drag-and-drop, the Gantt year timeline, and the Calendar weekly
        time-grid were all implemented and tested on at least two operating systems
        (Section~\ref{sec:res_views}).

  \item \textbf{Team collaboration and real-time communication.}
        Project and team invitations flow through MongoDB as a relay layer.
        Direct messages and group chats are persisted in SQLite and delivered in real time
        via a SignalR hub.
        File attachments up to 20~MB are supported.
        Notification delivery was verified for all seventeen event types
        (Sections~\ref{sec:res_teams} and~\ref{sec:res_settings}).

  \item \textbf{Integrated AI assistant.}
        The chatbot integrates the Mistral API in three modes: general productivity
        guidance, code generation and explanation, and document OCR and summarisation.
        Responses are streamed token-by-token over Server-Sent Events with a median
        first-token latency of approximately 400~ms in the general mode
        (Section~\ref{sec:res_ai}).

  \item \textbf{Maintainable, extensible codebase.}
        The backend follows a layered architecture (Controllers, Services, Repositories,
        Data) enforced by constructor injection via the ASP.NET Core DI container.
        The \texttt{ISyncableEntity} abstraction decouples the sync engine from
        individual entity types, allowing new syncable entities to be added by
        implementing the interface and registering a corresponding repository.
        AutoMapper profiles and FluentValidation validators are co-located with their
        respective DTOs and entities.

\end{enumerate}

% ─────────────────────────────────────────────────────────────────────────────
\section{Technical Contributions}\label{sec:concl_contributions}
% ─────────────────────────────────────────────────────────────────────────────

The project produced five technical contributions that go beyond assembling existing
libraries.

\paragraph{Outbox-pattern offline sync engine for a desktop .NET application.}
While the transactional outbox pattern is well documented in distributed-systems
literature, applying it within a single-user SQLite-backed desktop application required
adapting the pattern to a process-local context.
The \texttt{ConnectivityService} acts as the circuit breaker; the
\texttt{OfflineSyncService} provides the consumer loop; and the \texttt{ISyncableEntity}
interface defines the contract that entities must satisfy to participate in sync.
The result is a lightweight, dependency-free sync engine that requires no message broker
and adds less than one millisecond of overhead per write operation on the test machine.

\paragraph{Dual-database architecture with a shared entity abstraction.}
Combining an embedded relational database (SQLite) as the primary store with a
cloud-hosted document database (MongoDB) as a relay layer is an uncommon design for a
desktop application.
The \texttt{ISyncableEntity} interface exposes a \texttt{ToMongoDocument()} method that
produces a BSON-compatible dictionary, keeping the transformation logic on the entity
itself rather than scattered across service classes.
This design made it straightforward to add or modify syncable entities throughout
development without changing the sync engine.

\paragraph{SSE-streamed multi-modal AI chatbot in a desktop context.}
Integrating a streaming HTTP response (SSE) inside a Tauri-hosted application
required ensuring that the native WebView did not buffer or terminate the
long-lived HTTP connection before all tokens were delivered.
The solution --- issuing the SSE request directly from the React renderer process to
the local ASP.NET Core server --- bypassed the shell entirely and delivered
correct streaming behaviour.

\paragraph{Self-contained desktop packaging with embedded runtime.}
Bundling an ASP.NET Core server, a React frontend, and a Tauri shell into a single
installer that starts all three processes transparently on Windows
required careful orchestration of the Tauri main process.
The stdout-signal protocol (\texttt{TASKFLOW\_BACKEND\_READY:\{address\}}) decouples
the frontend load from the backend start time, allowing the UI to display a splash
state until the server is confirmed ready regardless of machine speed.

\paragraph{Tauri sidecar architecture with splashscreen lifecycle management.}
The Tauri v2~\cite{tauri_docs} desktop shell packages the ASP.NET Core backend
as a sidecar binary~\cite{tauri_sidecar}, launched and managed by Rust code
in \texttt{backend.rs}.
A two-window lifecycle (undecorated splashscreen during backend startup, then
main window with the full application) provides visual feedback without
requiring the user to understand the multi-process architecture.
The splashscreen closes atomically when the backend signals readiness via
stdout, ensuring the UI is never displayed before the API is available.
Nine IPC commands bridge the Rust shell and the web renderer, covering
backend URL discovery, startup status, password-reset code retrieval,
window management, and application metadata.
Neither the React~+~TypeScript frontend nor the ASP.NET Core
backend required any modification to accommodate the desktop shell,
confirming the framework-agnostic design of the three-tier architecture.
The resulting Tauri-based installer was approximately 35~MB in size with a
cold-start time of approximately 1.5~seconds on the test machine
(see Table~\ref{tab:perf_timings}).
% ─────────────────────────────────────────────────────────────────────────────
\section{Lessons Learned}\label{sec:concl_lessons}
% ─────────────────────────────────────────────────────────────────────────────

Several broader lessons emerged from the design and implementation process.

\textbf{Offline-first requires deliberate UI design, not just backend changes.}
It is not sufficient to enqueue writes; users need visible confirmation that their
actions have been recorded and will eventually reach the server.
The \texttt{ConnectivityBar} with its pending-count badge, colour-coded connectivity
states, and animated sync progress bar was as important to the user experience as the
outbox queue itself.
Implementing this required defining a clear contract between the sync service and the
React component so that both could react to the same state transitions.

\textbf{Choosing the right conflict resolution strategy early saves significant rework.}
The decision to use last-write-wins based on \texttt{UpdatedAt} was made early and
shaped the entity interface, the outbox entry schema, and the MongoDB upsert logic.
Had this decision been deferred, adding a timestamp field and revision counter to every
entity late in the project would have required a database migration on potentially live
data.

\textbf{SignalR group-per-user is a reliable notification pattern.}
Joining each authenticated connection to a group named after the user's ID
(\texttt{Groups.}\cbrk\texttt{AddToGroupAsync(}\cbrk\texttt{Context.ConnectionId, userId)}) meant that
notifications could be dispatched to any service class by calling
\texttt{\_hubContext.}\cbrk\texttt{Clients.}\cbrk\texttt{Group(userId)}, without needing to track connection IDs.
This pattern handled reconnection transparently: a reconnecting client re-joined its
group and received any notifications dispatched after reconnection without any
server-side state changes.

\textbf{Sidecar stdout parsing requires line-buffered input handling.}
The stdout-signal approach (\texttt{TASKFLOW\_BACKEND\_READY}) works well in
production builds but can be disrupted during development when the backend
writes diagnostic output before the ready signal.
Filtering the stdout stream with a prefix rather than reading the first line
proved more robust, but a dedicated named-pipe or socket
protocol would be a more principled solution.

\textbf{Clean architectural boundaries enable shell-layer substitution.}
Because the architecture separates the desktop shell from the web-based UI
and the process-based backend (Section~\ref{sec:meth_arch}), the Rust source
files in \texttt{src-tauri/} and the Tauri build pipeline are completely
decoupled from the React frontend and ASP.NET Core backend.
This validates the principle that investing in clean boundaries at the start of
a project pays dividends when technology choices need to be revisited.

% ─────────────────────────────────────────────────────────────────────────────
\section{Limitations}\label{sec:concl_limitations}
% ─────────────────────────────────────────────────────────────────────────────

The following limitations were identified during development and testing.

\textbf{Last-write-wins conflict resolution.}
The synchronisation engine uses a simple last-write-wins strategy based on the
\texttt{UpdatedAt} timestamp.
In scenarios where two users edit the same shared task concurrently while offline, the
later sync will silently overwrite the earlier one.
No merge dialogue or conflict notification is presented to the user.
For the single-user scenario (one user across multiple devices) this is acceptable;
for multi-user collaborative editing it may result in lost updates.

\textbf{Sub-task backend not implemented.}
The Table View renders a placeholder sub-task row beneath each task, and the task
entity schema includes a comment referring to sub-tasks.
However, no \texttt{SubTask} entity, migration, service, or API endpoint was created.
The feature exists at the UI layer only and is not persisted.
This is the largest gap between the stated scope and the delivered system.

\textbf{Mistral API key is hardcoded.}
The API key used to authenticate requests to the Mistral API is stored as a
compile-time constant in \texttt{MistralChatService.cs}.
This violates the principle of separating secrets from source code and means that
rotating the key requires a recompile and redistribution of the application.
The recommended remedy --- reading the key from the \texttt{TASKFLOW\_MISTRAL\_KEY}
environment variable at startup and falling back to a configuration value --- was
identified but not implemented within the project timeline.

\textbf{No automated end-to-end test suite.}
Testing was performed manually as described in Section~\ref{sec:meth_testing}.
While the architecture is amenable to unit testing (stateless services, injected
repositories), no unit tests, integration tests, or end-to-end tests were written.
The absence of an automated regression suite means that confidence in correctness relies
entirely on manual verification.

\textbf{macOS Gatekeeper requires a manual override on first launch.}
Because the macOS distribution was not submitted to Apple for notarisation, users on
macOS are required to right-click and select \emph{Open} on the first launch.
This is a distribution-process limitation rather than a functional defect but
represents a friction point for non-technical users.

%\let\cleardoublepage\clearpage
